\documentclass[11pt]{article}
\usepackage[margin=1in]{geometry}
\usepackage{amsmath,amssymb,booktabs,longtable,tabularx,array}
\usepackage{enumitem}
\usepackage[T1]{fontenc}
\usepackage[hidelinks]{hyperref}
\setlist[itemize]{leftmargin=*}
\setlist[enumerate]{leftmargin=*}
\newcommand{\code}[1]{\texttt{\detokenize{#1}}}
\newcommand{\st}{s_t}
\newcommand{\snext}{s_{t+1}}
\newcommand{\obs}{o_t}
\newcommand{\sobs}{\sigma_{\mathrm{obs}}}
\newcommand{\ssafe}{s_{\mathrm{safe}}}
\newcommand{\E}{\mathbb{E}}
\newcommand{\Var}{\mathrm{Var}}
\newcommand{\1}{\mathbf{1}}
\sloppy


\title{Algorithm Adaptations and Claims}
\author{Ecology Benchmark Documentation}
\date{July 2026}

\begin{document}
\maketitle

These are linear/mechanistic re-implementations of algorithmic ideas, not the
published deep-RL systems. The package depends on NumPy and PyYAML, and
\code{src/} contains no \code{torch}, \code{jax}, or \code{d3rlpy}. Dynamics
learning uses ridge-regularised linear ensembles; value fitting uses ridge
least-squares linear Q-functions; planning is finite-action particle MPC. The
comparison is therefore about control strategies under a shared ecological model
class, not about reproducing the original neural architectures.

\section{Identifier Mapping}

\begin{center}
\begin{tabularx}{\linewidth}{llll}
\toprule
Code identifier & Reader-facing name & Paper family & Model class here \\
\midrule
\code{mopo} & MOPO-style & Yu et al. 2020 MOPO & ridge-linear ensemble plus pessimistic particle MPC \\
\code{refplan} & RefPlan-inspired & Reflect-then-Plan & posterior ensemble planner \\
\code{bamcts} & BA-MCTS-inspired & BA-MCTS / BAMCP family & finite-action belief tree \\
\code{plus} & PLUS & BioConserv18 PLUS & mechanistic Ricker bank over candidate $K$ \\
\code{moor} & MOOR & fishery MOOR family & single misspecified Ricker, fits $K$ in set-point mode \\
\code{delphic} & Delphic-CQL-inspired & Delphic offline RL & linear/ridge compatible-world adaptation \\
\code{ogsrl} & OGSRL-inspired & Offline Guarded Safe RL & linear actor, kNN guardian, learned linear dynamics \\
\bottomrule
\end{tabularx}
\end{center}

The paper-family cells above were checked against the local method notes and
reference filenames, but final publication should still reconcile them with the
manuscript bibliography.

\section{Shared Infrastructure}

Most methods consume a public belief cache built from observations and public
controls. Learned transition models are bootstrap ridge members over
log-abundance, action one-hot terms, action/state interactions, and public
control features. Several methods share \code{ParticleMPC}, which samples finite
action sequences, rolls belief particles forward under a method-specific
proposal, and scores return with pessimism or risk terms.

\section{MOPO-Style}

\paragraph{Original idea.}
Learn a transition model from offline data, penalize uncertain model rollouts,
and avoid actions whose model predictions are unreliable.

\paragraph{Original architecture vs this implementation.}
The published MOPO family is typically neural model-based offline RL with
probabilistic ensembles and policy optimization. This implementation uses a
ridge-linear bootstrap ensemble and particle MPC; no neural actor or SAC-style
optimization is present.

\paragraph{What this code implements.}
\code{mopo.py} fits \code{ContinuousDynamicsEnsemble} and plans with
ensemble-disagreement and return-dispersion pessimism.

\paragraph{Adapted or simplified.}
Synthetic rollout policy optimization is replaced by receding-horizon finite
action planning to keep the benchmark deterministic, small, and inspectable.

\paragraph{Information access.}
It may use public observations, actions, logged rewards, public controls, and
belief features. It may not use private truth sidecars at training or
deployment.

\paragraph{Evaluation-time action.}
It chooses the first action of the best particle-MPC sequence.

\paragraph{Claim boundary.}
Valid claim: a MOPO-style pessimistic model-based planner in the benchmark's
linear model class. Invalid claim: reproduction of the published neural MOPO
system.

\section{RefPlan-Inspired}

\paragraph{Original idea.}
Maintain posterior uncertainty over dynamics hypotheses and plan under that
posterior.

\paragraph{Original architecture vs this implementation.}
The paper family uses richer Bayesian latent dynamics and trajectory-posterior
machinery. This implementation uses the same ridge-linear ensemble family as
MOPO and a residual observation-likelihood update over members.

\paragraph{What this code implements.}
\code{refplan.py} scores shared candidate action sequences under an ensemble
subset and subtracts posterior return uncertainty.

\paragraph{Adapted or simplified.}
The learned latent encoder and paper-specific planning machinery are replaced by
a transparent ensemble posterior and shared particle MPC.

\paragraph{Information access.}
Public dataset and public belief state only.

\paragraph{Evaluation-time action.}
Posterior-weighted sequence planning; execute the first action of the best
sequence.

\paragraph{Claim boundary.}
Valid claim: posterior-ensemble planning. Invalid claim: full doubly Bayesian
RefPlan architecture.

\section{BA-MCTS-Inspired}

\paragraph{Original idea.}
Plan in a belief over states and models with tree search, using exploration
bonuses and pessimistic model values.

\paragraph{Original architecture vs this implementation.}
The paper family includes continuous Bayes-adaptive search and learned policy or
value components. This implementation performs decision-time finite-action tree
search with bucketed continuous states and a ridge ensemble posterior.

\paragraph{What this code implements.}
\code{bamcts.py} root-samples belief particles and model members, applies UCB in
a finite action tree, and aggregates continuous states more finely near the
safety threshold.

\paragraph{Adapted or simplified.}
No paper-specific policy distillation loop is implemented; tree leaves and model
posterior are benchmark approximations.

\paragraph{Information access.}
Public belief and learned model ensemble only.

\paragraph{Evaluation-time action.}
Run tree search at the current belief and execute the selected root action.

\paragraph{Claim boundary.}
Valid claim: BA-MCTS-inspired online belief-tree planner. Invalid claim:
paper-faithful offline BA-MCTS training loop.

\section{PLUS}

\paragraph{Original idea.}
Maintain a finite bank of mechanistic candidate models, update candidate
evidence from observations, and act under the posterior mixture.

\paragraph{Original architecture vs this implementation.}
PLUS is already mechanistic, so the model-class gap is smaller than for the
neural offline-RL methods. In this set-point benchmark, however, the candidate
axis changes: the code spans candidate $K$ values because action-specific
growth set-points are known from the action table.

\paragraph{What this code implements.}
\code{plus.py} uses 21 Ricker-style candidates. Under set-point mode they are
geometrically spaced from $K_{\mathrm{base}}$ to $K_{\max}$ and each candidate
has its own belief-bank copy.

\paragraph{Adapted or simplified.}
The candidate prior is deliberately narrow: it does not include structural
family, process noise, or observation-noise candidate axes.

\paragraph{Information access.}
Public observations and controls; no true family parameters outside the
candidate model definition.

\paragraph{Evaluation-time action.}
Posterior-weighted candidate scores choose the first action of a sequence.

\paragraph{Claim boundary.}
Valid claim: close-to-intent mechanistic PLUS adaptation for known set-point
growth. Invalid claim: full PLUS candidate-model prior from the source paper.

\section{MOOR}

\paragraph{Original idea.}
Use a simple mechanistic ecological model as an interpretable expert baseline,
then plan under that fitted model.

\paragraph{Original architecture vs this implementation.}
MOOR is also mechanistic, so it is close to the intended spirit. The original
fishery pipeline is not reproduced; the code fits a Ricker model to public
belief means and plans with particle MPC over the benchmark action set.

\paragraph{What this code implements.}
\code{moor.py} fits a single Ricker proposal. Under set-point mode, $r$ is known
from $\rho$, so the fit identifies $K$ while mirroring the simulator's
positive-growth/negative-mortality convention.

\paragraph{Adapted or simplified.}
The model remains Ricker even when the true benchmark family is Allee, theta, or
regime-switching. That misspecification is intentional.

\paragraph{Information access.}
Public belief means, public controls, and logged rewards.

\paragraph{Evaluation-time action.}
Particle MPC under the fitted Ricker proposal.

\paragraph{Claim boundary.}
Valid claim: close-to-intent mechanistic Ricker expert baseline. Invalid claim:
full reproduction of the original MOOR fishery workflow.

\section{Delphic-CQL-Inspired}

\paragraph{Original idea.}
Penalize actions whose values vary across compatible latent explanations of the
same public evidence, then fit a conservative value function.

\paragraph{Original architecture vs this implementation.}
The paper family uses richer latent world models and neural/variational
components. This implementation uses random-feature latent worlds, ridge/linear
heads, and diagnostics over behavior-NLL and Bellman-error ranges.

\paragraph{What this code implements.}
\code{delphic.py} creates multiple public-feature worlds, estimates
counterfactual value variation, and subtracts a Delphic uncertainty penalty in a
linear CQL-style fitted-Q head.

\paragraph{Adapted or simplified.}
The worlds are benchmark-compatible approximations, not a proof that every
world induces the same observational distribution.

\paragraph{Information access.}
Public belief features and logged public rewards.

\paragraph{Evaluation-time action.}
Greedy action under linear Q minus Delphic penalty.

\paragraph{Claim boundary.}
Valid claim: Delphic-CQL-inspired conservative value adaptation. Invalid claim:
paper-faithful compatible-world causal implementation.

\section{OGSRL-Inspired}

\paragraph{Original idea.}
Learn an offline policy while guarding against unsafe and out-of-distribution
decisions.

\paragraph{Original architecture vs this implementation.}
The paper family may use neural policies, learned dynamics, and scalable support
estimators. This implementation uses a learned ridge-linear dynamics ensemble, a
kNN support guardian, a linear softmax actor, and linear diagnostic critics.

\paragraph{What this code implements.}
\code{ogsrl.py} trains with reward, safety, and support costs, updates dual
variables, and uses a deployment feasibility mask with a least-violating
fallback if every action violates a guard.

\paragraph{Adapted or simplified.}
The guardian and actor are intentionally low-capacity and benchmark-native.

\paragraph{Information access.}
Public belief features, public controls, and logged rewards. Safety estimates at
deployment are computed over posterior particles, not true hidden state.

\paragraph{Evaluation-time action.}
Select the highest-probability feasible action; if none is feasible, select the
least-violating action.

\paragraph{Claim boundary.}
Valid claim: OGSRL-inspired guarded safe offline control. Invalid claim:
transfer of the original theorem to these learned belief features.

\section{Summary}

PLUS and MOOR are closest to original intent because their source families are
already mechanistic. The general offline-RL methods are more distant: they carry
the strategy-level ideas into an auditable linear/mechanistic benchmark class.
That distinction is a strength for inspection, but it must be stated whenever
results are compared to the method names.

\section*{Verification}

Verified against: \path{pyproject.toml}, \path{src/real_ecology_benchmark/dynamics.py}, \path{src/real_ecology_benchmark/planning.py}, \path{src/real_ecology_benchmark/methods/value.py}, \path{src/real_ecology_benchmark/methods/mopo.py}, \path{src/real_ecology_benchmark/methods/refplan.py}, \path{src/real_ecology_benchmark/methods/bamcts.py}, \path{src/real_ecology_benchmark/methods/plus.py}, \path{src/real_ecology_benchmark/methods/moor.py}, \path{src/real_ecology_benchmark/methods/delphic.py}, \path{src/real_ecology_benchmark/methods/ogsrl.py}, \path{docs/real_ecology_history/29_6_algorithm_method_notes.tex}, \path{docs/real_ecology_history/AUDIT_algorithm_paper_implementation.md} @ be96c36

\end{document}
